\chapter{Epithelial Cells (Urine)}

\section*{What Is This Test?}
Epithelial cells in urine are exfoliated cells from the urinary tract lining. Squamous epithelial cells originate from the distal urethra and external genitalia; transitional (urothelial) cells from the renal pelvis, ureters, bladder, and proximal urethra; and renal tubular cells from the nephron. The clinical significance depends on cell type and quantity. Squamous cells are usually a contaminant from improper collection; transitional and especially renal tubular cells can signal pathology. Epithelial cell enumeration is part of the routine urine microscopic examination.

\section*{Alternative Names}
Urine Epithelial Cells; Urinalysis Epithelial Cells; Renal Tubular Cells; Transitional Cells; Squamous Cells (Urinary)

\section*{Principle}
Microscopic sediment examination at low (100x) and high power (400x) after centrifugation \cite{mcpherson2022henrys}. Squamous cells: large, flat, irregular borders, abundant cytoplasm, small central nucleus. Transitional cells: smaller, rounded or pear-shaped, larger nucleus. Renal tubular cells: rounded, slightly larger than WBCs, eccentric nucleus. Quantitated as few, moderate, many, or per high-power field.

\section*{History}
Microscopic urinary sediment examination has been performed since the 19th century. Modern automated imagen analyzers (IQ200, Atellica UAS800) classify cells using neural-network imaging; however, manual microscopy remains the reference standard.

\section*{Why This Test Is Ordered}
\begin{itemize}
  \item Evaluation of urinary tract infection (overwhelming squamous cells suggests contamination)
  \item Investigation of hematuria, proteinuria, or acute kidney injury
  \item Suspected acute tubular necrosis (renal tubular cells and muddy casts)
  \item Suspicion of viral infection (cytomegalovirus, BK virus) where viral inclusion-bearing tubular cells may be seen
  \item Routine urinalysis screening
\end{itemize}

\section*{Is This a Primary or Secondary Test}
Primary test, part of routine urinalysis.

\section*{How This Test Is Performed}
A 10--15 mL midstream clean-catch urine sample is centrifuged. The sediment is resuspended in 0.5 mL of supernatant, placed on a slide, and examined microscopically \cite{strasinger2020urinalysis}. A clean-catch technique reduces squamous cell contamination; catheter specimens are more reliable.

\section*{What Preparation Is Needed}
Clean-catch technique: cleanse perineum with provided antiseptic wipe, spread labia or retract foreskin, discard initial flow, collect midstream in sterile container. Avoid collection from the first morning tissue if menstrual blood is present.

\section*{Contraindications and Precautions}
Specimens $>$ 2 hours old at room temperature may show degenerated cells and increased bacteria. Refrigerated specimens can develop crystals that interfere with microscopy. Specimens contaminated with talc or lotion can show artifacts resembling cells.

\section*{Risks}
None.

\section*{What Happens After the Test}
Report integrated with dipstick findings. Many squamous cells suggest contamination; consider repeat specimen. Many renal tubular cells with casts suggest acute tubular necrosis \cite{perazella2015urine}.

\section*{Understanding Results}

\subsection*{Below Normal}
Few or no epithelial cells are normal. Absence of cells with strong clinical suspicion of tubular injury does not exclude acute tubular necrosis; tubular casts are more specific.

\subsection*{Above Normal}
\begin{itemize}
  \item \textbf{Many squamous cells:} Suggests contamination from external genitalia; specimen is unreliable for culture interpretation
  \item \textbf{Increased transitional cells:} Catheterization, recent instrumentation, bladder inflammation, urothelial tumors (atypical cells warrant cytology)
  \item \textbf{Renal tubular cells:} Pathologic; acute tubular necrosis (ischemic, nephrotoxic), acute interstitial nephritis, glomerulonephritis, allograft rejection (post-transplant), viral tubular infection (decoy cells in BK virus nephropathy)
  \item \textbf{Atypical urothelial cells:} Suspect urothelial carcinoma; requires urine cytology and cystoscopy
\end{itemize}

\section*{Normal Results}
\begin{itemize}
  \item \textbf{Squamous epithelial cells:} 0--2 per HPF (rare/occasional)
  \item \textbf{Transitional epithelial cells:} 0--1 per HPF
  \item \textbf{Renal tubular cells:} Absent or rare
  \item \textbf{Clinical context:} Many squamous cells require repeat clean-catch specimen
\end{itemize}

\section*{Follow-up Tests}
\begin{itemize}
  \item \textbf{Repeat urinalysis with clean-catch or catheter specimen:} If squamous contamination suspected
  \item \textbf{Urine culture:} If UTI suspected
  \item \textbf{Urine cytology:} If atypical urothelial cells or hematuria workup
  \item \textbf{Renal function tests (creatinine, BUN):} If renal tubular cells suggest acute kidney injury
  \item \textbf{Cystoscopy and upper-tract imaging:} In evaluation of painless hematuria \cite{aua2020microscopic}
  \item \textbf{Urine protein/creatinine ratio and microscopic casts:} For glomerular disease
\end{itemize}

\section*{References}
\bibliographystyle{unsrtnat}
\bibliography{references}

\clearpage